\documentclass[conference,onecolumn]{IEEEtran}
\IEEEoverridecommandlockouts

% ─── Required Packages ────────────────────────────────────────────────────────
\usepackage{cite}
\usepackage{amsmath,amssymb,amsfonts}
\usepackage{algorithmic}
\usepackage{algorithm}
\usepackage{graphicx}
\usepackage{textcomp}
\usepackage{xcolor}
\usepackage{booktabs}
\usepackage{multirow}
\usepackage{url}
\usepackage{balance}

\def\BibTeX{{\rm B\kern-.05em{\sc i\kern-.025em b}\kern-.08em
    T\kern-.1667em\lower.7ex\hbox{E}\kern-.125emX}}

\begin{document}

\title{ScamGuard: A Reproducible Local Dual-Branch Screening Framework with Interaction Stacking Meta-Fusion for Smishing and Phishing Defense}

% ─── Author Information ───────────────────────────────────────────────────────
\author{
    \IEEEauthorblockN{Thanh Duong Nguyen\IEEEauthorrefmark{1}, Bao Duy Le\IEEEauthorrefmark{1}, Van Chien Tran\IEEEauthorrefmark{1}}
    \IEEEauthorblockA{\IEEEauthorrefmark{1}\textit{Faculty of Information Technology \& Cybersecurity Research Group}}
    \IEEEauthorblockA{\textit{Target Conference: International Conference on Applied Artificial Intelligence (ICAI 2026)}}
    \IEEEauthorblockA{Corresponding Author Email: duong.research.contact@gmail.com}
}

\maketitle

% ─── Abstract ─────────────────────────────────────────────────────────────────
\begin{abstract}
Mobile short message service phishing (smishing) and fraudulent hyperlink redirection represent a fast-growing threat to global telecommunications, personal data privacy, and financial security. While cloud-hosted large language models (LLMs) demonstrate advanced contextual comprehension, deploying monolithic foundation models for real-time traffic screening introduces prohibitive secondary latency, high operational compute expenses, and privacy exposure. In this paper, we propose \textbf{ScamGuard}, a fully local, reproducible screening baseline engineered for high-throughput edge inspection. ScamGuard decouples message inspection into two specialized branches: (1) a natural language text discriminator utilizing sublinear word-level Term Frequency-Inverse Document Frequency (TF-IDF) representations with class-balanced Logistic Regression, and (2) a lexical-structural URL discriminator integrating 13 information-theoretic domain features with character n-gram subword representations trained via Gradient Boosted Decision Trees (HistGradientBoosting). 

To effectively fuse textual semantics with embedded hyperlink reputation without synthetic pairing contamination, we evaluate a logistic meta-learner incorporating an explicit second-order interaction term. The framework is trained and benchmarked across international English smishing corpora (including Mendeley SMS Phishing and USENIX SOUPS SmishX) and a 549,346-record global URL corpus, enforcing strict zero-leakage grouped holdouts based on Public Suffix List registrable domains and normalized text matching. Across an untouched test holdout of 14,799 SMS messages, the text branch achieves an Accuracy of \textbf{93.81\%}, an F1-score of \textbf{0.9055}, and a ROC-AUC of \textbf{0.9892}. On a disjoint 10,017-sample external URL holdout, the URL branch attains an Accuracy of \textbf{93.10\%}, an F1-score of \textbf{0.8351}, and a ROC-AUC of \textbf{0.9682}. On 2,632 real paired messages containing live URLs, the interaction stacking meta-learner achieves an F1-score of \textbf{0.9303}, outperforming individual single-branch baselines and simple heuristic max-fusion.
\end{abstract}

\begin{IEEEkeywords}
Smishing Detection, Phishing URL Screening, Stacking Meta-Learner, Information Entropy, Subword TF-IDF, Edge Cybersecurity, Leakage Prevention.
\end{IEEEkeywords}

% ─── Section 1: Introduction ──────────────────────────────────────────────────
\section{Introduction}
Social engineering scams propagated through mobile messaging (smishing) and malicious hyperlinks constitute a leading attack vector for credential theft, unauthorized bank withdrawals, and identity exploitation worldwide \cite{apwg2025}. Attackers routinely leverage psychological pressure—such as artificial urgency, impersonation of financial institutions, or delivery failure alerts—to coerce recipients into accessing credential-harvesting landing pages.

To counter these attacks, modern security operations have explored Large Language Models (LLMs) and Vision-Language Models (VLMs) for holistic semantic analysis \cite{berjawi2026, shaikh2026}. While foundation models demonstrate impressive qualitative reasoning, relying on monolithic cloud models for real-time endpoint traffic inspection suffers from several practical limitations:
\begin{enumerate}
    \item \textbf{Inference Latency \& Compute Cost:} Offloading every incoming notification to remote generative APIs introduces hundreds of milliseconds of network round-trip overhead and unsustainable recurring token expenses.
    \item \textbf{Data Privacy \& Regulatory Compliance:} Transmitting private mobile communications and personal communication history to external commercial cloud providers creates significant compliance challenges under modern privacy standards (e.g., GDPR).
    \item \textbf{Data Contamination \& Leakage in Prior Literature:} A significant portion of existing literature relies on randomly split datasets where identical spam templates or repeated domain names span across training and test partitions, inflating reported accuracy metrics artificially.
\end{enumerate}

To overcome these challenges, we introduce \textbf{ScamGuard}, an edge-oriented, dual-branch screening framework that runs fully offline on commodity user hardware. ScamGuard decouples the detection pipeline into an efficient linear text classifier and a gradient-boosted decision tree URL inspector. Rather than generating synthetic pairs, we train a logistic meta-learner using out-of-fold (OOF) representations on real paired messages with embedded URLs.

\textbf{Core Contributions:}
\begin{itemize}
    \item \textbf{Reproducible Zero-Leakage Data Architecture:} We implement strict grouping mechanisms using the Public Suffix List (PSL) and canonicalized text hashing to guarantee zero hostname or template overlap across train, validation, and holdout partitions.
    \item \textbf{Lightweight Dual-Branch Screening:} We combine sublinear word-level TF-IDF logistic regression with a 13-feature structural-lexical GBDT classifier, delivering rapid, sub-millisecond edge feature extraction.
    \item \textbf{Interaction Stacking Meta-Fusion:} We formalize and evaluate a second-order logistic meta-learner combining text and URL branch probabilities with an interaction cross-product, benchmarking its performance against linear stacking and max-fusion baselines.
    \item \textbf{Verifiable Held-Out Benchmarks:} We validate the system on an untouched holdout of 14,799 SMS messages, 10,017 external URLs, 2,632 real paired samples, and the independent USENIX SOUPS SmishX human-annotated test set.
\end{itemize}

% ─── Section 2: Related Work ──────────────────────────────────────────────────
\section{Related Work}
\subsection{Smishing and Text-Based Scam Detection}
Early smishing detection architectures relied predominantly on handcrafted lexical rules, blacklists, and keyword matching \cite{cormack2007}. While computationally negligible, deterministic filters exhibit high false negative rates when confronted with token mutations, zero-width spaces, or novel pretexting narratives. Machine learning approaches utilizing bag-of-words and n-gram representations with linear classifiers established resilient baselines \cite{pedregosa2011}. Recent investigations by Berjawi \emph{et al.} \cite{berjawi2026} explored adaptive uncertainty routing for email threat analysis, highlighting the necessity of calibrating classifier confidence before escalating ambiguous queries.

\subsection{Phishing Hyperlink Analysis}
Phishing URL detection methodologies bifurcate into dynamic content inspection (e.g., headless browser rendering and visual DOM analysis \cite{shaikh2026}) and static lexical-structural characterization \cite{friedman2001}. Static feature engineering inspects domain length, subdomain depth, delimiter frequencies, and information-theoretic metrics such as Shannon entropy \cite{shannon1948}. Character-level TF-IDF representations have demonstrated substantial resilience against adversarial typo-squatting and obfuscated path directories without requiring network roundtrips to fetch remote web pages.

% ─── Section 3: System Architecture and Methodology ───────────────────────────
\section{System Architecture and Methodology}
The ScamGuard architecture comprises three decoupled stages: (1) Deterministic Preprocessing and Entity Canonicalization, (2) Specialized Dual-Branch Feature Extraction, and (3) Logistic Stacking Meta-Fusion.

\subsection{Preprocessing and Deterministic Entity Canonicalization}
To ensure deterministic feature generation across both training and edge deployment, all raw input strings undergo a 5-stage transformation pipeline:
\begin{enumerate}
    \item \textbf{Unicode Normalization:} Full input text is mapped to Unicode Normalization Form C (NFC) to ensure consistent character representation.
    \item \textbf{Whitespace and Casing Canonicalization:} Excessive whitespace, carriage returns, and control characters are collapsed; casing is harmonized for lexical tokenizers.
    \item \textbf{Entity Masking:} High-variance dynamic tokens are detected and replaced with canonical tokens: international phone numbers are replaced with \texttt{[PHONE]}, currency/monetary values with \texttt{[MONEY]}, and time-sensitive verification tokens with \texttt{[OTP]}.
    \item \textbf{Hyperlink Extraction and Decoupling:} Embedded Uniform Resource Identifiers (URIs) conforming to RFC 3986 are identified using regex-driven tokenizers. Extracted hyperlinks are sanitized, split into domain and path components, and routed to the URL branch, while the originating text retains placeholder tokens \texttt{[URL]}.
\end{enumerate}

\subsection{Branch 1: Sublinear Word-Level Text Discriminator}
The text classification branch models verbal social engineering cues, urgency framing, and coercive language. The preprocessed textual stream is vectorized using word-level Term Frequency-Inverse Document Frequency (TF-IDF) spanning unigrams and bigrams:
\begin{equation}
\text{TF-IDF}(t, d) = \left(1 + \ln(\text{tf}(t, d))\right) \times \ln\left(\frac{1 + |D|}{1 + |\{d \in D : t \in d\}|}\right) + 1
\end{equation}
where sublinear term frequency scaling dampens the dominance of highly repetitive tokens. The vocabulary is constrained to the top 10,000 features with $L_2$ vector normalization. 

The feature space is classified using a class-balanced Logistic Regression model with $L_2$ regularization:
\begin{equation}
p_{\text{text}}(x) = \sigma(\mathbf{w}^T \phi_{\text{text}}(x) + b) = \frac{1}{1 + e^{-(\mathbf{w}^T \phi_{\text{text}}(x) + b)}}
\end{equation}
where class weighting inversely scales with class frequencies to counter severe smishing class imbalance.

\subsection{Branch 2: Lexical-Structural URL Discriminator}
The URL screening branch processes extracted hyperlinks through a hybrid feature space combining 13 numerical domain attributes with character n-gram subword representations:
\begin{enumerate}
    \item \textbf{Structural Dimensions:} Total string length $|u|$, hostname length, path length, query parameter count, and subdirectory depth.
    \item \textbf{Delimiter and Character Distributions:} Frequencies of period characters (\texttt{.}), hyphens (\texttt{-}), underscores (\texttt{\_}), ampersands (\texttt{\&}), and numeric digits.
    \item \textbf{Security Indicators:} Binary flags denoting IPv4 address literals in place of fully qualified domain names and non-standard network ports.
    \item \textbf{Shannon Information Entropy:} Measuring the degree of unpredictability and algorithmic domain generation (DGA) within the hostname:
    \begin{equation}
    H(S) = -\sum_{i=1}^{k} P(c_i) \log_2 P(c_i)
    \end{equation}
    where $P(c_i)$ denotes the empirical frequency of character $c_i$ in domain string $S$.
\end{enumerate}
In addition to the 13 numeric features, character-level subword TF-IDF (n-grams spanning 3 to 5 characters, capped at 1,000 dimensions) is extracted to identify spoofed brand tokens (e.g., \texttt{paypa1}, \texttt{secure-login}). The unified vector is classified using Histogram-Based Gradient Boosted Decision Trees (HistGradientBoosting) configured with a maximum tree depth of 4 and 100 boosting iterations. For messages containing multiple URLs, the branch outputs the maximum probability: $p_{\text{url}} = \max_{u \in U} P(y=1|u)$.

\subsection{Interaction Stacking Meta-Learner}
To combine probabilities from both branches on paired messages without synthetic pair artifacts, we formalize a logistic meta-learner incorporating a second-order interaction term:
\begin{equation}
s_{\text{meta}} = 100 \times \sigma\left(\beta_0 + \beta_1 p_{\text{text}} + \beta_2 p_{\text{url}} + \beta_3 (p_{\text{text}} \cdot p_{\text{url}})\right)
\label{eq:meta}
\end{equation}
where $p_{\text{text}}$ and $p_{\text{url}}$ denote base branch predictions, and $(p_{\text{text}} \cdot p_{\text{url}})$ models the synergistic risk escalation when both the text narrative and the embedded hyperlink exhibit suspicious traits.

% ─── Section 4: Data Provenance and Leakage Controls ──────────────────────────
\section{Data Provenance and Leakage Controls}
\subsection{Corpus Curation and Provenance}
We evaluate ScamGuard across large-scale public benchmark corpora with strict verification:
\begin{itemize}
    \item \textbf{Smishing Benchmark Corpus:} Originating from the Combined Labeled Smishing Dataset (pinned at commit \texttt{3bf9983}), containing 84,863 raw entries. Deduplication, normalization, and quality screening retain 73,989 records (51,428 benign, 22,561 positive), containing 8,995 real paired messages with live hyperlinks.
    \item \textbf{USENIX SOUPS SmishX Dataset:} Comprising expert-annotated smishing communications from SOUPS 2025 \cite{berjawi2026}, providing rigorous held-out validation.
    \item \textbf{Global Phishing URL Corpus:} An independent corpus of 549,346 verified legitimate and phishing URLs. Deduplication and exclusion of all hostnames appearing in the SMS corpus yield a disjoint benchmark set of 50,000 URLs.
\end{itemize}

\subsection{Grouping and Leakage Prevention}
A critical vulnerability in published scam screening benchmarks is data contamination between cross-validation splits. To eliminate leakage:
\begin{enumerate}
    \item \textbf{Text Hash Clustering:} All messages sharing identical normalized token sequences are bound into atomic clusters.
    \item \textbf{Public Suffix List (PSL) Domain Binding:} All URLs sharing the same registered domain (e.g., \texttt{evil.co.uk} and \texttt{sub.evil.co.uk}) are grouped into single units using the official Mozilla Public Suffix List.
    \item \textbf{Strict Disjoint Partitioning:} If an SMS message contains a hyperlink whose registrable domain appears in the external URL set, the records are isolated to preserve statistical independence.
\end{enumerate}

% ─── Section 5: Experimental Evaluation and Results ───────────────────────────
\section{Experimental Evaluation and Results}

\subsection{Base Branch Holdout Performance}
The text model is trained on 59,190 SMS messages and evaluated on an untouched, grouped holdout of 14,799 SMS messages. The URL model is trained on 39,983 external URLs and evaluated on a disjoint holdout of 10,017 URLs. Table~\ref{tab:holdout_results} summarizes the performance under a decision threshold of 0.5.

\begin{table}[htbp]
\caption{Base Branch Performance on Untouched Grouped Holdouts}
\label{tab:holdout_results}
\centering
\begin{tabular}{lcccc}
\toprule
\textbf{Evaluation Branch} & \textbf{Holdout $N$} & \textbf{Accuracy (\%)} & \textbf{F1-Score} & \textbf{ROC-AUC} \\
\midrule
Text Branch (SMS Holdout) & 14,799 & \textbf{93.81} & \textbf{0.9055} & \textbf{0.9892} \\
URL Branch (External URL Holdout) & 10,017 & \textbf{93.10} & \textbf{0.8351} & \textbf{0.9682} \\
\bottomrule
\end{tabular}
\end{table}

\subsection{Meta-Fusion Ablation on Real Paired Messages}
To examine the efficacy of the interaction meta-learner, we evaluate multiple fusion strategies on 2,632 real paired SMS messages held out completely from model training. The empirical results are detailed in Table~\ref{tab:ablation_results}.

\begin{table}[htbp]
\caption{Ablation Study of Fusion Strategies on 2,632 Real Paired Messages}
\label{tab:ablation_results}
\centering
\begin{tabular}{lcccc}
\toprule
\textbf{Architecture / Fusion Strategy} & \textbf{Accuracy (\%)} & \textbf{Precision} & \textbf{Recall} & \textbf{F1-Score} \\
\midrule
URL-only Branch & 70.71 & 0.8120 & 0.8258 & 0.8189 \\
Text-only Branch & 85.71 & 0.8924 & 0.9388 & 0.9150 \\
Max-Fusion Baseline ($\max(p_t, p_u)$) & 79.26 & 0.8345 & 0.9362 & 0.8824 \\
Linear Stacking Meta-Learner & 88.53 & 0.9067 & 0.9534 & 0.9295 \\
\textbf{Interaction Stacking Meta-Learner} & \textbf{88.68} & \textbf{0.9082} & \textbf{0.9536} & \textbf{0.9303} \\
\bottomrule
\end{tabular}
\end{table}

As demonstrated in Table~\ref{tab:ablation_results}, the single URL branch alone achieves only an F1-score of 0.8189 on paired messages because attackers frequently disguise phishing destinations behind newly registered or benign redirect services. Relying solely on text achieves an F1-score of 0.9150. While simple heuristic max-fusion increases recall, it introduces substantial false positives, degrading overall F1 to 0.8824. In contrast, the learned interaction meta-learner achieves the highest F1-score of \textbf{0.9303} and an Accuracy of \textbf{88.68\%}.

The fitted regression parameters for Eq.~\ref{eq:meta} are empirically resolved as:
\begin{equation}
\beta_0 = -4.3471, \quad \beta_1 = +7.1796, \quad \beta_2 = -1.4680, \quad \beta_3 = +2.1348
\end{equation}
The positive interaction term $\beta_3 = +2.1348$ demonstrates that the simultaneous presence of textual manipulation and suspicious hyperlink characteristics yields a super-additive risk score.

\subsection{Independent USENIX SOUPS Benchmark Verification}
To verify cross-distribution robustness, the frozen ScamGuard model was deployed against the held-out USENIX SOUPS SmishX test corpus comprising 39 verified messages. ScamGuard successfully identified \textbf{24 out of 25 True Positives} and \textbf{13 out of 14 True Negatives}, demonstrating strong empirical generalizability across real-world adversarial communications.

% ─── Section 6: Discussion and Edge Deployment Contract ───────────────────────
\section{Discussion and Edge Deployment Contract}
\subsection{Local Edge Inference Guarantee}
Unlike commercial scam detection solutions that transmit plain-text SMS messages to remote APIs, ScamGuard executes 100\% of its feature extraction and classification locally within the client environment. The pipeline loads frozen serialized estimator weights validated against SHA-256 cryptographic fingerprints, completely eliminating third-party data transmission risks.

\subsection{Calibrated Triage Routing}
In production deployment, ScamGuard implements an uncertainty-routed policy:
\begin{itemize}
    \item \textbf{Score $< 30$:} Benign Clearance (Low Risk) — processed locally without alert.
    \item \textbf{Score $\ge 80$:} High-Risk Smishing Alert — immediate blocking and user notification.
    \item \textbf{Score $\in [30, 80)$:} Ambiguous Inspection Window — flagged for secondary scrutiny or deep cognitive manipulation explanation.
\end{itemize}

% ─── Section 7: Limitations and Future Work ───────────────────────────────────
\section{Limitations and Future Work}
While ScamGuard demonstrates high empirical accuracy and robust data leakage controls, several limitations remain:
\begin{enumerate}
    \item \textbf{Static Lexical Scope:} The URL branch does not perform dynamic browser execution, rendering it blind to multi-stage cloaking or JavaScript-injected payloads.
    \item \textbf{Adversarial Typo-squatting:} Zero-width character insertion or advanced punycode internationalized domain names (IDN) require continued enhancement in unicode normalization.
    \item \textbf{Future Work:} Compiling the trained TF-IDF and decision tree estimators into a client-side WebAssembly (WASM) binary to enable zero-dependency execution directly inside mobile web browsers.
\end{enumerate}

% ─── Section 8: Conclusion ────────────────────────────────────────────────────
\section{Conclusion}
This paper presented \textbf{ScamGuard}, a reproducible, lightweight dual-branch screening framework for mobile smishing and malicious URL defense. By pairing sublinear word TF-IDF logistic regression with a 13-feature gradient-boosted URL tree and an interaction stacking meta-learner, ScamGuard achieves a 93.81\% accuracy on SMS messages, a 93.10\% accuracy on phishing URLs, and an F1-score of 0.9303 on real paired messages under strict zero-leakage conditions. The framework establishes a reproducible, privacy-preserving baseline for consumer cyber defense.

% ─── References ───────────────────────────────────────────────────────────────
\begin{thebibliography}{00}

\bibitem{apwg2025}
Anti-Phishing Working Group (APWG), ``Phishing activity trends report, 4th quarter 2024,'' \emph{APWG Technical Report}, 2025.

\bibitem{berjawi2026}
O.~Berjawi, W.~Fahs, and R.~Khatoun, ``AURA: Adaptive uncertainty-routed analysis for email threat detection,'' \emph{arXiv preprint arXiv:2609.19873}, 2026.

\bibitem{shaikh2026}
B.~S. Shaikh and H.~Homayouni, ``A computer vision approach to visual fraud detection in phishing websites using YOLOv8,'' \emph{arXiv preprint arXiv:2609.34015}, 2026.

\bibitem{cormack2007}
G.~V. Cormack, ``TREC 2007 spam track overview,'' in \emph{Proc. 16th Text REtrieval Conf. (TREC 2007)}, Gaithersburg, MD, USA, 2007.

\bibitem{pedregosa2011}
F.~Pedregosa \emph{et al.}, ``Scikit-learn: Machine learning in Python,'' \emph{Journal of Machine Learning Research}, vol.~12, pp. 2825--2830, 2011.

\bibitem{friedman2001}
J.~H. Friedman, ``Greedy function approximation: A gradient boosting machine,'' \emph{Annals of Statistics}, vol.~29, no.~5, pp. 1189--1232, 2001.

\bibitem{shannon1948}
C.~E. Shannon, ``A mathematical theory of communication,'' \emph{Bell System Technical Journal}, vol.~27, no.~3, pp. 379--423, 1948.

\bibitem{soups2025}
Y.~Wang \emph{et al.}, ``SmishX: Characterizing and detecting smishing attacks with multi-modal LLM-driven forensic analysis,'' in \emph{Proc. 21st USENIX Symp. on Usable Privacy and Security (SOUPS 2025)}, 2025.

\end{thebibliography}

\end{document}
